\chapter{配列と繰り返し}
\section{配列の作成と利用} \label{sec:02-array}
\begin{itembox}[l]{配列}
\begin{verbatim}
 1: #include <FK/FK.h>
 2: 
 3: int main(int, char **)
 4: {
 5:     fk_AppWindow    window;
 6:     fk_Block        block(10.0, 10.0, 10.0);
 7:     fk_Model        model[2];
 8: 
 9:     model[0].setShape(&block);
10:     model[0].setMaterial(Yellow);
11:     model[0].glMoveTo(-20.0, -20.0, 0.0);
12:     window.entry(&model[0]);
13: 
14:     model[1].setShape(&block);
15:     model[1].setMaterial(Red);
16:     model[1].glMoveTo(20.0, 20.0, 0.0);
17:     window.entry(&model[1]);
18: 
19:     window.setSize(800, 600);
20:     window.setBGColor(0.3, 0.6, 0.8);
21:     window.open();
22: 
23:     while(window.update() == true) {
24:     }
25: 
26:     return 0;
27: }
\end{verbatim}
\end{itembox}
\subsection*{解説}
\begin{itemize}
 \item 7行目で、これまでの変数の準備(これを「変数の\textbf{宣言}」といいます)
	の際に、カギ括弧で 2 という数字を挟んでいます。
	このように、変数宣言の際にカギ括弧で数字を挟んでいる形式の
	変数を「\textbf{配列}」と言います。
	配列は、沢山の変数をひとまとまりで準備することができます。
	今回のように「\verb+model[2]+」とした場合は、
	「\verb+model[0]+」と「\verb+model[1]+」という変数を
	準備したことと同じ意味となります。このように、
	「\verb+配列名[n]+」と宣言した場合、
	\verb+配列名[0]+から\verb+配列名[n-1]+ の n 個の変数が
	利用できるようになります。

 \item 9行目と14行目では、同じ block を異なるモデルに設定しています。
	このように、FK では異なるモデルに同じ形状を設定することが可能です。
	例えば、カーレースゲームでは同じ形状のモデルに対して
	色を替えた車を複数同時に扱いたいものです。
	このようなとき、FK では各車のモデルを用意し、
	それぞれのモデルに同じ形状を設定するということができます。
	モデルはそれぞれ固有の色や位置を持つことができます。
	今回も、色や位置が二つのモデルで異なっていることがわかります。

\end{itemize}

\section{for文によるループ} \label{sec:02-for}
\begin{itembox}[l]{for文}
\begin{verbatim}
 1: #include <FK/FK.h>
 2: 
 3: int main(int, char **)
 4: {
 5:     fk_AppWindow    window;
 6:     fk_Block        block(1.0, 1.0, 1.0);
 7:     fk_Model        model[10];
 8:     int             i;
 9: 
10:     for(i = 0; i < 10; i++) {
11:         model[i].setShape(&block);
12:         model[i].setMaterial(Yellow);
13:         model[i].glMoveTo(i*5, 0.0, 0.0);
14:         window.entry(&model[i]);
15:     }
16: 
17:     window.setSize(800, 600);
18:     window.setBGColor(0.3, 0.6, 0.8);
19:     window.open();
20: 
21:     while(window.update() == true) {
22:     }
23: 
24:     return 0;
25: }
\end{verbatim}
\end{itembox}
\subsection*{解説}
\begin{itemize}
 \item 8行目、int という型の「\verb+i+」という変数が宣言されています。
	この変数については、これまでの変数と違って
	「情報を格納する」という目的で変数を利用していきます。
	int 型は、整数の値を1個だけ格納することができます。

 \item 10行目の「for」というのは変数ではなく、
	C++の中で「ループ」を扱うための機能です。
	書式は以下の様になっています。\\ ~ \\
	\begin{shadebox}
	\begin{verbatim}
	for(初期処理; 繰り返し条件; 繰り返し処理) {
	        実行文1;
	        実行文2;
	           :
	}
	\end{verbatim}
	\end{shadebox} ~ \\
	for文では、以下の様な処理が行われます。
	\begin{enumerate}
	\item まず、「初期処理」を実行する。今回の場合は
		i に 0 が代入される。
	\item 次に、「繰り返し条件」で条件判定が行われる。ここでは、
		i が 10 未満かどうかを判定する。もしこの判定に
		合格した場合、「\verb+{+」と「\verb+}+」で囲まれた
		「実行文」が逐次実行されていく。
	\item 全ての実行文が実行された時点で、「繰り返し処理」
		が実行される。今回の場合は「\verb-i++;-」と
		書いてあるが、これは「i に 1 を足す」という意味である。
		例えば、i に 0 が入っていたときは 1 に、
		3 が入っていたときは 4 になる。
	\item 後は、ステップ 2 に戻って同様のことが行われていく。
		もし判定に合格しなかった場合は、「\verb+}+」の
		外を実行していく。
	\end{enumerate}
	結果として、この for 文の中身は 10 回実行され、
	また実行の度に i が 1 ずつ足されていくということになっています。

 \item 11〜13行目では、model 配列のカギ括弧の中の数字
	(これを\textbf{添字}と言います)が、直接の数値指定ではなく
	変数 i を用いています。このように、
	添字には変数や変数の計算結果(例えば\verb-i+3- や \verb+i*2+ など)
	を入れることが可能です。

 \item 13行目の glMoveTo 関数では、第一引数として i に 5 倍した数値を
	入力しています。結果的に、最初の箱は \((0, 0, 0)\)、
	二個目の箱は \((5, 0, 0)\)、三個目の箱は \((10, 0, 0)\)
	というように、それぞれ別の場所に移動することになります。

\end{itemize}
\newpage
\section{練習課題} \label{sec:02-q}
\begin{description}
 \myitem \ref{sec:02-for}節のサンプルプログラムに対し、
	10個の箱がウィンドウの真ん中を中心に並ぶように修正せよ。

 \myitem 10個の箱が斜めに並ぶように配置せよ。

 \myitem 100個の箱が \(10\times10\) のマス目状に並ぶように配置せよ。

 \myitem マス目状に並んでいる100個の箱に対し、
	箱の色が赤と黄色で市松模様(チェック模様)になるようにせよ。
	ただし、if文はまだ用いないこと。
	(ヒント: 配列は複数宣言することが可能)
\end{description}
\begin{figure}[H]
\centering
\includegraphics[scale=0.4]{./Fig/Fig02-01.eps}
\end{figure}
